% Generated draft: requires mathematical and source review.
\documentclass[11pt]{article}
\usepackage[margin=25mm]{geometry}
\usepackage{fontspec}
\setmainfont{Latin Modern Roman}
\usepackage{amsmath,amssymb,mathtools,unicode-math}
\setmathfont{Latin Modern Math}
\usepackage{enumitem,needspace,xurl,hyperref}
\setlength{\parindent}{0pt}
\setlength{\parskip}{5pt}
\setlength{\emergencystretch}{3em}
\newcommand{\N}{\mathbb{N}}
\newcommand{\Z}{\mathbb{Z}}
\newcommand{\Q}{\mathbb{Q}}
\newcommand{\R}{\mathbb{R}}
\newcommand{\C}{\mathbb{C}}
\newcommand{\F}{\mathbb{F}}
\newcommand{\A}{\mathbb{A}}
\newcommand{\PP}{\mathbb{P}}
\newcommand{\E}{\mathbb{E}}
\newcommand{\eps}{\varepsilon}
\newcommand{\dd}{\,\mathrm{d}}
\begin{document}
\section*{UnsolvedMath ID: 30001825}
\textbf{Source problem code:} catalogue entry
\subsection{Definitions and mathematical statement}
Let \((M,g)\) be a complete four-dimensional Riemannian manifold. Write \(\operatorname{Rm}_g\) for its Riemann curvature tensor and \(\mathrm{dvol}_g\) for its Riemannian volume measure. The finite curvature energy condition is
\[
\operatorname{Rm}_g\in L^2(M,g)
\quad\Longleftrightarrow\quad
\int_M |\operatorname{Rm}_g|_g^2\,\mathrm{dvol}_g<\infty.
\]
Consider the class
\[
\mathcal C=\left\{(M,g):\begin{array}{l}
\dim M=4,\quad g\text{ is complete},\\
(M,g)\text{ is a gravitational instanton in the source's sense},\\
\displaystyle\int_M|\operatorname{Rm}_g|_g^2\,\mathrm{dvol}_g<\infty,\\
(M,g)\text{ is neither ALE nor ALF}
\end{array}\right\}.
\]
Here ALE and ALF denote asymptotically locally Euclidean and asymptotically locally flat, respectively. The supplied record does not specify the geometric models, decay estimates, or end conditions defining these predicates. It also does not specify the convention for gravitational instantons, including whether hyperk\"ahler structure is required or a broader Ricci-flat class is intended. Consequently, \(\mathcal C\) is a conditional description of the intended class, pending these conventions.

\textbf{Classification problem.} Classify all members of \(\mathcal C\), with the equivalence relation for the classification to be specified. Determine whether members can have slow curvature decay. More specifically, the original formulation asks whether the geometries described as newly described in the source exhaust the members with slow curvature decay. Determine also whether such metrics can occur as singularity dilations.

\textbf{Unresolved qualifications.} The supplied record neither identifies those newly described geometries nor defines slow curvature decay. It does not specify the ambient geometric problem, singularities, rescaling procedure, or convergence notion meant by singularity dilations. The referent of such metrics in the last question could be all members of \(\mathcal C\), its slow-decay members, or the newly described examples. These missing specifications prevent a fully quantified classification, exhaustiveness, or singularity-limit statement; no choice among them is made here.
\subsection{Short English statement}
Classify complete four-dimensional gravitational instantons of finite total squared curvature that are neither ALE nor ALF. The source also asks about slow curvature decay, whether its unnamed new geometries exhaust that behavior, and realization as singularity dilations. Definitions, the candidate geometries, and the intended singularity setting require clarification before these questions become precise.
\subsection{Sources}
\begin{itemize}
\item Local UnsolvedMath registry, supplied record with permanent ID \(30001825\), title \emph{Classification of Finite-Energy Gravitational Instantons}, source problem code \texttt{catalogue entry}; collection: Oberwolfach Reports Open Problems. The proposed year is recorded as 2011; no individual proposer is supplied.
\item Supplied report citation: \emph{Differentialgeometrie im Gro{\ss}en (2012)}, Oberwolfach Reports, DOI 10.4171/owr/2011/33, pp.\ 1857--1912. Original Oberwolfach report: \url{https://doi.org/10.4171/owr/2011/33}. The supplied background describes these reports as records of mathematical workshops and their open-problem discussions.
\item Supplied literature reference labeled ``Li classification results'': \url{https://arxiv.org/abs/2310.13197}. No title, precise theorem, or detailed hypotheses for these results are supplied.
\item Literature assessment dated exactly \texttt{2026-08-21}: the supplied curation labels the problem partially solved and reports that recent classification results cover substantial conformally K\"ahler and Hermitian subfamilies with finite curvature energy, while not establishing the broad classification conjecture described in the extracted record. This is supplied dated curation, not an independently verified literature assessment; its time-sensitive status and cited claims require independent verification.
\end{itemize}
\end{document}
